\subsection{Canonicalization Preserves Feasibility}
\label{app:canonicalization-proof}
We prove Lemma~\ref{lem:canonicalization} under the assumptions of
Section~\ref{sec:formal-scope}, with every input axis having extent at least
two, as required by Lemma~\ref{lem:axis-exactness}. We use the retained-axis
and contraction-role rules in Appendix~\ref{app:canonical-axis-rules}.
Start with a valid attribute assignment
$\beta$. By Lemma~\ref{lem:axis-exactness}, its axes can be identified
with the LHS logical axes.

\subsubsection{Proof Plan}
We construct $\beta'$ from the inputs toward the root. At each node $v$,
we show that $P_v$ is a subset of the original output axes and that legal
attributes can produce exactly
\begin{equation}
 \operatorname{Ord}_v(\beta')=\operatorname{Ord}_v(\beta)|_{P_v}.
\end{equation}
In words, we delete axes outside $P_v$ without changing the order of those
that remain. Every mandatory axis must also survive under $\beta$:
once removed, it cannot connect to another input or reach the final output.

\subsubsection{Constructing Attributes at Each Node}
\paragraph{Inputs and unary operations.}
Inputs stay unchanged. A permutation can arrange its retained axes in
the required order. A unary summation keeps $P_v=P_u\cap M_v$ and sums
out the others. These mandatory axes also survive under $\beta$, and
summation preserves their order.

\paragraph{Contraction.}
Assume the construction holds for both children $\ell,r$.
There are three points to check:
\begin{enumerate}
\item \emph{Shared axes must connect here.}
An axis in $P_\ell\cap P_r$ occurs in both original children.
Under $\beta$, these occurrences must be paired at this contraction;
otherwise they become separate axes in one subtree and cannot be joined
later.
\item \emph{Shared axes have the required roles.}
An axis in $B_v$ is mandatory, so $\beta$ must keep it as batch.
An axis in $K_v$ has no connections outside the subtree and is absent
from the final output. We can contract it now, even if $\beta$ keeps
it as batch and sums it out later.
D2 permits this change of axis roles and the corresponding changes to
position lists.
\item \emph{Free axes remain free.}
Suppose a left-free axis $a\in P_\ell\setminus P_r$ appeared in the
original right child. Since it connects inputs in both children,
it would be mandatory in the right child and hence belong to $P_r$,
a contradiction. Thus it was left-free under $\beta$ too.
The right-free case is identical.
\end{enumerate}
Therefore every axis in $P_v=B_v\cup LF_v\cup RF_v$ survives in the
original output with the same role. Keep the original relative order
of batch axes, pair contraction axes in the same order on both sides,
and inherit the retained free-axis orders from the children.
This produces exactly the original output order restricted to $P_v$.
The attributes are legal: positions are distinct, batch and contraction
axes are disjoint, and paired axes have equal extents.

\subsubsection{Why the Result Is Valid}
The construction connects only matching LHS logical axes. It removes an
axis only after all its input connections are complete and only if the
axis is absent from the final output. Thus the input-port connections
are unchanged.

At the root, the mandatory axes and the original output axes are both
$\operatorname{set}(F)$. The subset property above gives
\begin{equation}
 \operatorname{set}(F)\subseteq P_{\mathrm{root}}
 \subseteq\operatorname{set}(F).
\end{equation}
Hence the root keeps exactly those axes, in the original order $F$.
The new assignment is legal and has the same axis representation as
$\beta$, so Lemma~\ref{lem:axis-exactness} makes it valid.
